\worldchapter{\trans{quirks_title}}{appendix-quirks}
\label{ch:appendix-quirks}

\begin{multicols}{2}

    \section{Wahnsinn und Visionen}

    \subsection*{Schizophrenie}

    Der Charakter erlebt die Welt auf eine Weise, die von der Realität anderer Menschen abweicht. Er hat möglicherweise Halluzinationen, durch die er Dinge sieht, hört oder fühlt, die nicht vorhanden sind. Zudem kann er Wahnvorstellungen haben, die sein Handeln und Denken beeinflussen.

    Der Charakter kann außergewöhnlich kreativ sein und einzigartige Lösungen für Probleme finden, auf die andere nicht kommen würden. Seine andere Art, die Welt zu sehen, kann zu innovativen Ideen und künstlerischen Ausdrucksformen führen. Er verfügt über eine ausgeprägte \textbf{Darbietungs-} und \textbf{Logikfähigkeit}.

    Die Symptome der Schizophrenie erschweren es dem Charakter, soziale Kontakte zu knüpfen und stabile Beziehungen zu führen. \textbf{Kommunikationsfähigkeit} und \textbf{Charme} sind eingeschränkt.

    \begin{tabular}{@{}ll@{}}
        \textbf{Logik} & +1\\
        \textbf{Charme} & -1\\
        \textbf{Darbietung} & +2\\
        \textbf{Kommunikation} & -2\\
    \end{tabular}

    \subsection*{Wahn}

    Der Charakter hält an Überzeugungen fest, die nicht der Realität entsprechen und von anderen nicht geteilt werden. Diese Wahnvorstellungen können verschiedene Themen betreffen, beispielsweise Verfolgungswahn, Größenwahn oder Beziehungswahn.

    Der Wahn verleiht dem Charakter ein starkes Gefühl von Zweck und Richtung. Er glaubt unerschütterlich an seine Sache, was ihm in bestimmten Situationen eine starke innere Motivation und Entschlossenheit verleiht. Seine \textbf{Tapferkeit} und \textit{Willenskraft}* sind erhöht.

    Wahnvorstellungen können den Charakter von der Realität entfremden und seine Fähigkeit beeinträchtigen, rational zu handeln oder effektiv mit anderen zu interagieren. Der Spielleiter kann jederzeit eine Logik-Probe verlangen. Misslingt diese, hat der Charakter in diesem Moment eine aktive Wahnvorstellung.

    \begin{tabular}{@{}ll@{}}
        \textbf{Willenskraft} & +1\\
        \textbf{Logik} & +1\\
    \end{tabular}

    \subsection*{Dunkle Visionen}

    Der Charakter wird von unheimlichen Visionen heimgesucht, die ihm gelegentlich Einblicke in verborgene oder zukünftige Ereignisse gewähren. Diese Visionen sind häufig fragmentarisch und rätselhaft. Manchmal zeigen sie beunruhigende oder bedrohliche Szenen. Sie kommen unvorhersehbar und lassen den Charakter manchmal wie in einer anderen Realität gefangen erscheinen.

    Die Visionen können dem Charakter wertvolle Hinweise oder Warnungen liefern, die ihm dabei helfen, Gefahren zu vermeiden und verborgene Wahrheiten aufzudecken.

    Die Visionen sind oft beunruhigend und können den Charakter emotional belasten. Nach einer Vision muss der Charakter einen Stresstest ablegen. Misslingt dieser, erhält er 2 \textbf{Stress}.


    \subsection*{Besessenheit}

    Der Charakter wird von einer fremden Entität beeinflusst, die ihm zwar zusätzliche Fähigkeiten verleiht, aber auch gelegentlich die Kontrolle über seinen Körper oder Geist übernimmt. Diese Besessenheit manifestiert sich in unvorhersehbaren Momenten und kann sowohl nützlich als auch gefährlich sein. Die Entität verfolgt eigene Ziele und Absichten, die nicht immer mit denen des Charakters übereinstimmen.

    Der Charakter erhält einen permanenten Bonus von +2 auf ein Attribut seiner Wahl, das der Entität entspricht.

    Gelegentlich kann die Entität die Kontrolle übernehmen und den Charakter zu unvorhersehbaren oder gefährlichen Handlungen zwingen. In dem Moment, in dem die Entität den Charakter übernimmt, wirft dieser einen W6. Je höher das Ergebnis, desto positiver ist das Eingreifen für den Charakter. Ergebnisse über 3 sind in der Regel positiv.


    \subsection*{Zweifelnder Geist}

    Der Charakter ist von Natur aus misstrauisch und hinterfragt ständig die Absichten und Handlungen anderer. Diese Skepsis schützt ihn zwar vor Täuschungen, macht es ihm aber auch schwer, echte Verbündete zu finden und Vertrauen aufzubauen. Er neigt dazu, in allem einen versteckten Plan oder eine Falle zu sehen.

    Der Charakter ist immun gegen Manipulationen und Täuschungen, die einen Wurf auf Täuschen und Betrügen involvieren.

    Aufgrund seines ständigen Misstrauens hat der Charakter Schwierigkeiten, anderen zu vertrauen. Er erhält -2 auf seinen Wert in \textbf{Kommunikation}.

    \begin{tabular}{@{}ll@{}}
        \textbf{Kommunikation} & -2\\
    \end{tabular}

    \section{Impulsives Verhalten}

    \subsection*{Depression}

    Der Charakter leidet unter anhaltenden Gefühlen von Hoffnungslosigkeit. Er hat möglicherweise das Interesse an Aktivitäten verloren, die ihm früher Freude bereitet haben, und fühlt sich oft müde oder energielos.

    Aufgrund seiner Erfahrungen mit Depressionen kann der Charakter ein tiefes Verständnis und Mitgefühl für das Leid anderer entwickeln.

    Die Symptome der Depression können den Charakter daran hindern, alltägliche Aufgaben zu bewältigen oder sich an sozialen Aktivitäten zu beteiligen. Er kann sich zurückziehen und isolieren.

    \begin{tabular}{@{}ll@{}}
        \textbf{Gewissenhaftigkeit} & -1\\
        \textbf{Ausdauer} & -1\\
        \textbf{Empathie} & +2\\
    \end{tabular}

    \subsection*{Gedankenlos}

    Der Charakter handelt oft impulsiv und ohne lange nachzudenken. Diese Spontanität führt dazu, dass er schnell Entscheidungen trifft und sich auf seine Instinkte verlässt. Allerdings kann diese Eigenschaft auch dazu führen, dass ihm wichtige Details entgehen oder er unüberlegte Risiken eingeht.

    Der Charakter ist in der Lage, in gefährlichen oder stressigen Situationen schnell zu handeln. Bei der Regel der \textbf{Blitzreaktion} kann er einen fehlgeschlagenen Wurf einmal wiederholen.

    Aufgrund seiner impulsiven Art neigt der Charakter dazu, wichtige Details zu übersehen und unüberlegte Risiken einzugehen. Dies hat einen Malus auf \textbf{Logik} zur Folge.

    \begin{tabular}{@{}ll@{}}
        \textbf{Logik} & -1\\
    \end{tabular}

    \subsection*{Rastlosigkeit}

    Der Charakter ist von innerer Unruhe getrieben, die ihn ständig in Bewegung hält. Er hat Schwierigkeiten, zur Ruhe zu kommen oder sich über einen längeren Zeitraum an einem Ort aufzuhalten.

    Aufgrund seiner Rastlosigkeit kann der Charakter länger ohne Pause aktiv bleiben. Seine Ausdauer ist erhöht.

    Der Charakter hat Probleme, sich auf Aufgaben zu konzentrieren, die Ruhe oder Geduld erfordern. \textbf{Geschick} und \textbf{Gewissenhaftigkeit} sind verringert.

    \begin{tabular}{@{}ll@{}}
        \textbf{Gewissenhaftigkeit} & -1\\
        \textbf{Geschick} & -1\\
        \textbf{Ausdauer} & +2\\
    \end{tabular}

    \subsection*{Draufgänger}

    Der Charakter wird zunehmend waghalsiger und fahrlässiger. Er geht Risiken ein, wo andere zögern würden, und handelt oft aus dem Bauch heraus.

    Der Charakter hat eine erhöhte \textbf{Tapferkeit}.

    Die \textbf{Logik} des Charakters ist verringert.

    \begin{tabular}{@{}ll@{}}
        \textbf{Logik} & -2\\
        \textbf{Tapferkeit} & +2\\
    \end{tabular}

    \section{Persönliche Eigenheiten}

    \subsection*{Nervosität}

    Der Charakter ist von grundlegender Nervosität geprägt. Diese sorgt einerseits dafür, dass er im Kampf leicht abgelenkt ist, ermöglicht ihm aber auch, überraschende Angriffe auszuführen.

    Pro Kampf erhält der Charakter 10 Würfel, die er beliebig auf Angriffe und Verteidigungen anwenden kann.

    Der Charakter kann keine kritischen Treffer im Angriff erzielen.


    \subsection*{Ungeschickt}

    Der Charakter neigt dazu, Dinge fallen zu lassen, zu stolpern oder auf andere Weise ungeschickt zu sein. Diese Eigenschaft kann im Alltag zu unbeabsichtigten Missgeschicken führen. Sie bringt jedoch auch eine gewisse Unvorhersehbarkeit und Kreativität mit sich.

    Der Charakter zeichnet sich durch seine einzigartige Fähigkeit aus, Probleme auf unkonventionelle Weise zu lösen. Seine ungeschickte Art führt oft zu Lösungen, die anderen nicht in den Sinn kommen würden und die daher unerwartet sind. Dafür erhält er einen Bonus von zwei Würfeln auf Proben, die unkonventionelles Denken erfordern.

    Der Charakter hat ein verringertes \textbf{Geschick}.

    \begin{tabular}{@{}ll@{}}
        \textbf{Geschick} & -2\\
    \end{tabular}

    \subsection*{Schwacher Lebenswille}

    Der Charakter zeigt eine geringe Widerstandsfähigkeit gegenüber den Herausforderungen des Lebens. In schwierigen Situationen neigt er dazu, schneller aufzugeben oder sich von Rückschlägen entmutigen zu lassen.

    Aufgrund seiner Gleichgültigkeit hat sich der Charakter in gewisser Hinsicht seinem Schicksal ergeben. Er erhält einen Schicksalswürfel.

    Der Charakter hat eine verringerte \textbf{Resistenz}.

    \begin{tabular}{@{}ll@{}}
        \textbf{Schicksalswürfel} & +1\\
        \textbf{Resistenz} & -2\\
    \end{tabular}

    \subsection*{Perfektionismus}

    Der Charakter strebt in allem, was er tut, nach Perfektion. Er gibt sich nie mit halbherzigen Lösungen zufrieden und arbeitet unermüdlich daran, seine Tätigkeiten bestmöglich zu erledigen.

    Aufgrund seines Perfektionismus erhält der Charakter zwei Bonuswürfel und zwei Wiederholungswürfel.

    Komplexere Tätigkeiten des Charakters dauern in der Regel doppelt so lange wie bei anderen Charakteren. Falls es auf Geschwindigkeit ankommt, kann der Spielleiter einen Logikwurf verlangen. Falls dieser misslingt, verdoppelt sich die angegebene Zeit.

    \begin{tabular}{@{}ll@{}}
        \textbf{Wiederholungswürfe} & +2\\
        \textbf{Bonuswürfel} & +2\\
    \end{tabular}

    \section{Ängste und Abneigungen}

    \subsection*{Agoraphobie}




    \begin{tabular}{@{}ll@{}}
        \textbf{Ausweichen} & +1\\
        \textbf{Gewissenhaftigkeit} & -1\\
        \textbf{Wahrnehmung} & +1\\
    \end{tabular}

    \subsection*{Soziale Phobie}




    \begin{tabular}{@{}ll@{}}
        \textbf{Wiederholungswürfe} & -1\\
        \textbf{Gewissenhaftigkeit} & +1\\
        \textbf{Kommunikation} & -1\\
    \end{tabular}

    \subsection*{Akrophobie}




    \begin{tabular}{@{}ll@{}}
        \textbf{Wahrnehmung} & +1\\
        \textbf{Tapferkeit} & -2\\
    \end{tabular}

    \subsection*{Hydrophobie}




    \begin{tabular}{@{}ll@{}}
        \textbf{Tapferkeit} & -2\\
        \textbf{Orientierung} & +1\\
    \end{tabular}

    \subsection*{Mysophobie}




    \begin{tabular}{@{}ll@{}}
        \textbf{Gewissenhaftigkeit} & +1\\
        \textbf{Resistenz} & -1\\
    \end{tabular}

    \subsection*{Zoophobie}




    \begin{tabular}{@{}ll@{}}
        \textbf{Ausweichen} & +1\\
        \textbf{Tapferkeit} & -1\\
        \textbf{Naturkunde} & -1\\
    \end{tabular}

    \subsection*{Dermatozoenwahn}




    \begin{tabular}{@{}ll@{}}
        \textbf{Maximaler Stress} & +1\\
        \textbf{Gewissenhaftigkeit} & -2\\
    \end{tabular}

    \subsection*{Coulrophobie}




    \begin{tabular}{@{}ll@{}}
        \textbf{Auffassungsgabe} & +1\\
        \textbf{Darbietung} & -1\\
        \textbf{Tapferkeit} & -1\\
    \end{tabular}

    \subsection*{Teratophobie}




    \begin{tabular}{@{}ll@{}}
        \textbf{Wiederholungswürfe} & +1\\
        \textbf{Resistenz} & -1\\
        \textbf{Orientierung} & -1\\
    \end{tabular}

    \subsection*{Verfolgungsangst}

    der betroffene wird von Phasen geplagt in denen er sich von einer paranoiden Verfolgungsangst heimgesucht sieht,



    \begin{tabular}{@{}ll@{}}
        \textbf{Charme} & -1\\
        \textbf{Heimlichkeit} & -1\\
        \textbf{Tapferkeit} & -1\\
    \end{tabular}

    \subsection*{Angst vor Altertümern}

    Der Charakter hat eine irrationale Angst vor allem Alten. Er ist überzeugt, dass in jedem alten Objekt eine Macht oder ein Geheimnis schlummert, das ihn bedroht.

    Aufgrund seiner Angst beschäftigt sich der Charakter natürgemäß mehr mit Altertümern. Er erhält den Wert 4 in Altertumskunde.

    In der Nähe von Altertümern (wie alten Gemäuern, Ruinen, Statuen, aber auch Schriftstücken oder Büchern) wird der Charakter durch die Angst ungeschickter. Alle \textbf{Physis} Attribute sind um eins verringert.

    \begin{tabular}{@{}ll@{}}
        \textbf{Altertümer} & +4\\
    \end{tabular}

    \subsection*{Ungläubig}

    Der Charakter hat jeglichen Glauben an höhere Wesen oder Religionen verloren. Er sucht für alles nach rationalen Erklärungen und ist oft misstrauisch gegenüber Dingen, die sich nicht logisch erklären lassen.

    Aufgrund seines Unglaubens ist der Charakter stets bemüht, eine rationale Erklärung zu finden. Er erhält drei \textbf{Wiederholungswürfe}.

    Alle Würfe, die Übernatürliches involvieren, haben einen um zwei erschwerten Mindestwurf. Der Charakter hat eine verminderte Religionsfertigkeit.

    \begin{tabular}{@{}ll@{}}
        \textbf{Religion} & -3\\
    \end{tabular}

    \subsection*{Hypochondrie}

    Der Charakter ist übermäßig besorgt um seine Gesundheit und deutet selbst die kleinsten körperlichen Symptome als Anzeichen für ernsthafte Krankheiten.

    Aufgrund seiner ständigen Sorge um seine Gesundheit ist der Charakter sehr gut über medizinische Themen und gesundheitliche Vorsorge informiert. Er erhält einen Bonus auf \textbf{Erste Hilfe} und das Wissen \textbf{Medizin}.

    Die übertriebene Angst vor Krankheiten belastet den Charakter stark und beeinträchtigt seine Handlungsfähigkeit. Seine \textbf{Logik} und \textbf{Gewissenhaftigkeit} sind verringert.

    \begin{tabular}{@{}ll@{}}
        \textbf{Logik} & -1\\
        \textbf{Gewissenhaftigkeit} & -1\\
        \textbf{Erste Hilfe} & +2\\
        \textbf{Medizin} & +1\\
    \end{tabular}

    \section{Stimmungsschwankungen}

    \subsection*{Nachtmensch}

    Eine Schlaflosigkeit in der Nacht plagt den Charakter und lässt ihn tagsüber erschöpft sein.

    Der Charakter ist nachts besonders wachsam und aufmerksam. Er erhält einen Bonus von +2 auf \textbf{Wahrnehmung} und +1 auf \textbf{Orientierung} bei Nacht.

    Der übermüdete Charakter erhält einen Malus von -2 auf \textbf{Logik} und von -1 auf \textbf{Mechanik} bei Tag.


    \subsection*{Launenhaft}

    Der Charakter ist für seine schnellen und unvorhersehbaren Stimmungsschwankungen bekannt. Diese Launenhaftigkeit lässt ihn manchmal charmant und überzeugend wirken, während er im nächsten Moment gereizt oder melancholisch erscheint.

    In sozialen Situationen kann der Charakter sehr charmant und überzeugend sein. Das verleiht ihm einen Bonus auf seinen \textbf{Charme}. Dank seiner Fähigkeit, schnell zwischen verschiedenen Stimmungen zu wechseln, kann er sich flexibel an unterschiedliche soziale Dynamiken anpassen.

    Bei jeder längeren oder tiefergehenden Unterhaltung muss der Charakter auf seine Kommunikation würfeln. Misslingt der Wurf, kommt es zu einer Stimmungsschwankung und der Bonus wird für die nächsten W6 Stunden zum Malus.

    \begin{tabular}{@{}ll@{}}
        \textbf{Charme} & +2\\
        \textbf{Kommunikation} & +1\\
    \end{tabular}

    \subsection*{Tagaktiv}

    Der Charakter ist tagsüber am leistungsfähigsten und fühlt sich bei Tageslicht besonders wohl. Nachts hingegen plagt ihn ein ständiges Unwohlsein und eine subtile Angst.

    Tagsüber ist der Mindestwurf für alle Fertigkeitswürfe um eins verringert.

    In der Nacht ist der Mindestwurf für alle Fertigkeitswürfe um eins erhöht.


    \subsection*{Rastloser Geist}

    Der Charakter hat einen rastlosen Geist, der ihn ständig wachsam und aufmerksam macht. Diese innere Unruhe verhindert, dass er jemals zur Ruhe kommt oder sich vollständig entspannen kann.

    Der Charakter ist stets aufmerksam und wachsam. Das verleiht ihm einen Bonus auf \textbf{Wahrnehmung}.

    Aufgrund seiner Unfähigkeit, vollständig zur Ruhe zu kommen, regeneriert der Charakter bei jeder Instanz von Stressreduktion nur die Hälfte der normalen Stressreduktion (aufgerundet). Auch bei der Rast wird der Stress nur halbiert statt auf 0 gesetzt.

    \begin{tabular}{@{}ll@{}}
        \textbf{Wahrnehmung} & +2\\
    \end{tabular}

    \section{Sinnesveränderungen}

    \subsection*{Hypersensivität}

    Der Charakter nimmt seine Umgebung intensiver wahr als andere. Er nimmt Geräusche, Lichter, Gerüche und sogar subtile Veränderungen in seiner Umgebung stärker und deutlicher wahr. Diese erhöhte Sensibilität kann sowohl eine Bereicherung als auch eine Belastung darstellen.

    Der Charakter erhält einen Bonus von +2 auf die Fertigkeit Wahrnehmung.

    In lauten Umgebung ist der Charakter schnell überfordert. In diesen Situationen muss er einen Stresstest bestehen. Misslingt dieser, so erhält er 1 \textbf{Stress}.

    \begin{tabular}{@{}ll@{}}
        \textbf{Wahrnehmung} & +2\\
    \end{tabular}

    \subsection*{Tinnitus}

    Der Charakter leidet unter einem konstanten Klingeln oder Summen in den Ohren, das ihn ständig begleitet. In ruhigen Momenten kann dieses Geräusch besonders störend sein und die Konzentration beeinträchtigen.

    Aufgrund der konstanten Geräusche in seinen Ohren ist der Charakter immun gegen akustische Täuschungen oder Manipulationen. Er kann nicht durch Geräusche aus seiner Umgebung abgelenkt werden.

    Der Mindestwurf aller \textbf{Wahrnehmungsproben}, die das Hören einschließen, ist um 3 erhöht.


    \section{Mystische Eigenschaften}

    \subsection*{Blutlinie der Alten}

    Der Charakter glaubt, dass er von einer uralten, mystischen Blutlinie abstammt. Dies soll ihm besondere Fähigkeiten sowie ein tiefes Verständnis für vergessene Sprachen und Symbole verleihen.

    Der Charakter kann alte Sprachen und Symbole verstehen und interpretieren. Er erhält einen Bonus von 2 Würfeln auf Proben auf Bildung, Historie und Altertümer, sofern diese im Zusammenhang mit alten Schriften stehen. Der Charakter erhält das Wissen Altertümer 1.

    Der Charakter fühlt sich elitär und wird zunehmend verschrobener. Er nimmt einen Malus auf Kommunikation und Charme in Kauf.

    \begin{tabular}{@{}ll@{}}
        \textbf{Charme} & -1\\
        \textbf{Kommunikation} & -1\\
        \textbf{Altertümer} & +1\\
    \end{tabular}

    \subsection*{Unstillbarer Wissensdurst}

    Der Charakter ist von einem unersättlichen Verlangen nach Wissen getrieben. Ständig ist er auf der Suche nach neuen Informationen, sei es durch Bücher, Gespräche oder Erkundungen.

    Der Charakter erhält einen Bonus auf \textbf{Logik} und \textbf{Bildung}.

    Der unstillbare Wissensdurst hat zur Folge, dass die generelle Wahrnehmung des Charakters leidet.

    \begin{tabular}{@{}ll@{}}
        \textbf{Bildung} & +1\\
        \textbf{Logik} & +1\\
        \textbf{Wahrnehmung} & -1\\
    \end{tabular}

    \subsection*{Asket}

    Der Charakter sieht zunehmend keinen Sinn mehr in materiellen Besitztümern und zieht sich auch sonst immer weiter zurück. Diese Lebensweise verleiht ihm zwar innere Stärke und Unabhängigkeit, führt aber auch dazu, dass er sich von sozialen Normen und Beziehungen distanziert.

    Der Charakter verspührt später Hunger und Durst, und er hat eine erhöhte \textbf{Resistenz}.

    Aufgrund seines enthaltsamen Lebensstils hat der Charakter Schwierigkeiten, sich in soziale Gruppen einzufügen oder Beziehungen aufzubauen. Dies führt zu einem Malus in \textit{Kommunikation}*.

    \begin{tabular}{@{}ll@{}}
        \textbf{Resistenz} & +2\\
        \textbf{Kommunikation} & -2\\
    \end{tabular}

    \section{Körperliche Besonderheiten}

    \subsection*{Schwächlich}

    Der Charakter wird zunehmend schwächer, seine Ängste lassen seine Kräfte schwinden. Einerseits kann ihn diese körperliche Schwäche anfälliger für physische Belastungen machen, andererseits verleiht sie ihm aber auch eine gewisse Leichtigkeit und Beweglichkeit.

    Aufgrund seiner weniger robusten Statur kann sich der Charakter leiser und unauffälliger bewegen. Er erhält einen Bonus auf \textbf{Heimlichkeit}.

    Der Charakter hat eine verringerte \textbf{Kraft} und ein ständiges Gefühl von körperlicher Schwäche und Unsicherheit.

    \begin{tabular}{@{}ll@{}}
        \textbf{Kraft} & -1\\
        \textbf{Heimlichkeit} & +2\\
    \end{tabular}

    \subsection*{Abhängigkeit}

    Der Charakter ist von einer bestimmten Substanz, Aktivität oder einem bestimmten Verhalten abhängig, um normal zu funktionieren. Diese Abhängigkeit kann sich auf verschiedene Lebensbereiche auswirken und die Entscheidungen und Handlungen massiv beeinflussen.

    Hat der Charakter Zugang zu dem, wovon er abhängig ist, ist er extrem fokussiert und motiviert. In dieser Zeit hat er zwei Schicksalswürfel.

    Ohne Zugang zu dem, wovon er abhängig ist, leidet der Charakter unter Entzugserscheinungen, welche seine körperliche und geistige Leistungsfähigkeit beeinträchtigen.  Sein Mindestwurf ist um eins erhöht.


\end{multicols}